% from ega4-2-fr.tex
\oldpage[IV-2]{215}

il suffit qu'il soit universellement caténaire et que ses fibres formelles
soient géométriquement régulières.

(ii) Si $A$ est un anneau excellent, il en est de même de tout anneau de
fractions $S^{-1}A$ et de toute $A$-algèbre de type fini.

(iii) Un anneau local complet (en particulier un corps) est excellent. Un
anneau de Dedekind dont le corps des fractions est de caractéristique $0$ (en
particulier $\mathbf Z$) est excellent.

(iv) Soient $A$ un anneau excellent, $X=\operatorname{Spec}(A)$; l'ensemble
$\operatorname{Reg}(X)$ (resp. $\operatorname{Nor}(X)$,
$\mathrm U_{\mathrm R_n}(X)$) des points où $X$ est régulier (resp. normal,
resp. vérifie $(\mathrm R_n)$) est ouvert dans $X$. Pour tout
$\mathcal O_X$-Module cohérent $\mathcal F$, l'ensemble
$\mathrm U_{\mathrm C_n}(\mathcal F)$ (resp.
$\mathrm U_{\mathrm S_n}(\mathcal F)$, $\operatorname{CM}(\mathcal F)$) des
$x\in X$ où $\operatorname{coprof}(\mathcal F_x)\leq n$ (resp. des points où
$\mathcal F$ vérifie $(\mathrm S_n)$, resp. des points où $\mathcal F_x$ est un
$\mathcal O_x$-Module de Cohen-Macaulay) est ouvert dans $X$.

(v) Soient $A$ un anneau excellent, $\mathfrak J$ un idéal de $A$,
$\widehat A$ le séparé complété de $A$ pour la topologie $\mathfrak J$-préadique.
Alors le morphisme canonique
$f:\operatorname{Spec}(\widehat A)\to\operatorname{Spec}(A)$ est régulier
(autrement dit, plat et à fibres géométriquement régulières). Si l'on pose
$X=\operatorname{Spec}(A)$, $X'=\operatorname{Spec}(\widehat A)$, on a (avec
les notations de (iv) et en posant
$\mathcal F'=\mathcal F\otimes_{\mathcal O_X}\mathcal O_{X'}$),
\begin{env}[7.8.3.1]
\label{IV.2.s7.8.3.1-fr}
\begin{align*}
\operatorname{Reg}(X')&=f^{-1}(\operatorname{Reg}(X)),&
\operatorname{Nor}(X')&=f^{-1}(\operatorname{Nor}(X)),&
\mathrm U_{\mathrm R_n}(X')&=f^{-1}(\mathrm U_{\mathrm R_n}(X)),\\
\mathrm U_{\mathrm C_n}(\mathcal F')&=f^{-1}(\mathrm U_{\mathrm C_n}(\mathcal F)),&
\mathrm U_{\mathrm S_n}(\mathcal F')&=f^{-1}(\mathrm U_{\mathrm S_n}(\mathcal F)),&
\operatorname{CM}(\mathcal F')&=f^{-1}(\operatorname{CM}(\mathcal F))
\end{align*}
\end{env}

En particulier, si $\mathfrak J$ est contenu dans le radical de $A$ (par
exemple si $A$ est local et $\mathfrak J$ son idéal maximal), pour que $A$ soit
régulier (resp. normal, resp. réduit, resp. vérifie $(\mathrm R_n)$, resp. soit
de coprofondeur $\leq n$, resp. vérifie $(\mathrm S_n)$, resp. soit un anneau
de Cohen-Macaulay), il faut et il suffit qu'il en soit de même de $\widehat A$;
en particulier, pour que $A$ n'ait pas d'idéaux premiers associés immergés, il
faut et il suffit qu'il en soit de même de $\widehat A$.

(vi) Un anneau excellent $A$ est universellement japonais; en particulier, si
$A$ est intègre, sa fermeture intégrale dans toute extension finie de son corps
des fractions est une $A$-algèbre finie.

(vii) Soit $A$ un anneau local excellent réduit. Alors le complété $\widehat A$
est réduit, la fermeture intégrale $A'$ de $A$ dans son anneau total des
fractions est une $A$-algèbre finie (donc un anneau semi-local), et la fermeture
intégrale de $\widehat A$ dans son anneau total des fractions est isomorphe au
complété $\widehat {A'}$ de $A'$. En outre, il y a une correspondance biunivoque
canonique entre l'ensemble des points maximaux de
$\operatorname{Spec}(\widehat A)$ (autrement dit, l'ensemble des idéaux premiers
minimaux de $\widehat A$) et l'ensemble des points fermés de
$\operatorname{Spec}(A')$ (autrement dit, l'ensemble des idéaux maximaux de
$A'$). Pour que l'anneau local $A$ soit unibranche, il faut et il suffit que
$\widehat A$ soit intègre; pour que $A$ soit géométriquement unibranche, il faut
et il suffit que $\widehat A$ le soit.

(viii) Si $A$ est un anneau excellent intègre, la clôture intégrale de $A$ est
l'intersection des clôtures intégrales des anneaux $A_{\mathfrak p}$, où
$\mathfrak p$ parcourt l'ensemble des idéaux premiers de $A$ de hauteur $1$.

(ix) Soient $A$ un anneau excellent, $X=\operatorname{Spec}(A)$, $Z$ une partie
fermée de $X$, $U=X-Z$, $i:U\to X$ l'injection canonique, $\mathcal F$ un
$\mathcal O_U$-Module cohérent. Pour que le $\mathcal O_X$-Module
$i_*(\mathcal F)$ soit cohérent, il faut et il suffit que, pour tout
$x\in\operatorname{Ass}(\mathcal F)$, on ait
$\operatorname{codim}(\overline{\{x\}}\cap Z,\overline{\{x\}})\geq2$. En
particulier, si $A$ est intègre et $\mathcal F$ sans torsion, pour que
$i_*(\mathcal F)$ soit cohérent, il faut et il suffit que
$\operatorname{codim}(Z,X)\geq2$.

